\documentclass[
	ngerman,
	toc=listof, % Abbildungsverzeichnis sowie Tabellenverzeichnis in das Inhaltsverzeichnis aufnehmen
	toc=bibliography, % Literaturverzeichnis in das Inhaltsverzeichnis aufnehmen
	footnotes=multiple, % Trennen von direkt aufeinander folgenden Fußnoten
	parskip=half, % vertikalen Abstand zwischen Absätzen verwenden anstatt horizontale Einrückung von Folgeabsätzen
	numbers=noendperiod % Den letzten Punkt nach einer Nummerierung entfernen (nach DIN 5008)
]{scrartcl}
\pdfminorversion=5 % erlaubt das Einfügen von pdf-Dateien bis Version 1.7, ohne eine Fehlermeldung zu werfen (keine Garantie für fehlerfreies Einbetten!)
\usepackage[utf8]{inputenc} % muss als erstes eingebunden werden, da Meta/Packages ggfs. Sonderzeichen enthalten
\usepackage[T1]{fontenc} % Wichtig für die Verwendung von T1-kodierten Schriftarten
\usepackage{uarial} % Paket laden, das die Schriftart bereitstellt
\renewcommand{\sfdefault}{ua1} % Setze die Standardschriftart auf Arial
\renewcommand{\familydefault}{\sfdefault} % Setze die Standardschriftfamilie auf serifenlos (Arial)


% !TEX root = Projektdokumentation.tex

% Hinweis: der Titel muss zum Inhalt des Projekts passen und den zentralen Inhalt des Projekts deutlich herausstellen
\newcommand{\titel}{Automatisierte Verfolgung und Benachrichtigung für verbesserte Liefertrans-parenz}
\newcommand{\untertitel}{Überwachungssystem für Warenlieferungen}
\newcommand{\kompletterTitel}{\titel{} -- \untertitel}

\newcommand{\autorName}{Jens Lange}
\newcommand{\autorAnschrift}{Hoßberg 17}
\newcommand{\autorOrt}{21376 Salzhausen}

\newcommand{\betriebLogo}{LogoBetrieb.pdf}
\newcommand{\betriebName}{\textsc{Vonmählen GmbH} }
\newcommand{\betriebAnschrift}{Vor dem Bardowicker Tore 49}
\newcommand{\betriebOrt}{21339 Lüneburg}

\newcommand{\ausbildungsberuf}{Fachinformatiker für Anwendungsentwicklung}
\newcommand{\betreff}{Dokumentation zur betrieblichen Projektarbeit}
\newcommand{\pruefungstermin}{Winter 2024}
\newcommand{\abgabeOrt}{Lüneburg}
\newcommand{\abgabeTermin}{27.11.2024}
 % Metadaten zu diesem Dokument (Autor usw.)
% !TEX root = ../Projektdokumentation.tex

% Anpassung an Landessprache ---------------------------------------------------
\usepackage{babel}
\usepackage{tablefootnote}
\usepackage[table,xcdraw]{xcolor} % For coloring table rows
\usepackage{array}               % For advanced table formatting

% Umlaute ----------------------------------------------------------------------
%   Umlaute/Sonderzeichen wie äüöß direkt im Quelltext verwenden (CodePage).
%   Erlaubt automatische Trennung von Worten mit Umlauten.
% ------------------------------------------------------------------------------
\usepackage[T1]{fontenc}
\usepackage{textcomp} % Euro-Zeichen etc.

% Schrift ----------------------------------------------------------------------
\usepackage{lmodern} % bessere Fonts
\usepackage{relsize} % Schriftgröße relativ festlegen

% Tabellen ---------------------------------------------------------------------
\PassOptionsToPackage{table}{xcolor}
\usepackage{tabularx}
% für lange Tabellen
\usepackage{longtable}
\usepackage{array}
\usepackage{ragged2e}
\usepackage{lscape}
\newcolumntype{w}[1]{>{\raggedleft\hspace{0pt}}p{#1}} % Spaltendefinition rechtsbündig mit definierter Breite

% Grafiken ---------------------------------------------------------------------
\usepackage[dvips,final]{graphicx} % Einbinden von JPG-Grafiken ermöglichen
\usepackage{graphics} % keepaspectratio
\usepackage{floatflt} % zum Umfließen von Bildern
\graphicspath{{Bilder/}} % hier liegen die Bilder des Dokuments

% Sonstiges --------------------------------------------------------------------
\usepackage[titles]{tocloft} % Inhaltsverzeichnis DIN 5008 gerecht einrücken
\usepackage{amsmath,amsfonts} % Befehle aus AMSTeX für mathematische Symbole
\usepackage{enumitem} % anpassbare Enumerates/Itemizes
\usepackage{xspace} % sorgt dafür, dass Leerzeichen hinter parameterlosen Makros nicht als Makroendezeichen interpretiert werden

\usepackage{makeidx} % für Index-Ausgabe mit \printindex
\usepackage[printonlyused]{acronym} % es werden nur benutzte Definitionen aufgelistet

% Einfache Definition der Zeilenabstände und Seitenränder etc.
\usepackage{setspace}
\usepackage{geometry}

% Symbolverzeichnis
\usepackage[intoc]{nomencl}
\let\abbrev\nomenclature
\renewcommand{\nomname}{Abkürzungsverzeichnis}
\setlength{\nomlabelwidth}{.25\hsize}
\renewcommand{\nomlabel}[1]{#1 \dotfill}
\setlength{\nomitemsep}{-\parsep}

\usepackage{varioref} % Elegantere Verweise. „auf der nächsten Seite“
\usepackage{url} % URL verlinken, lange URLs umbrechen etc.

\usepackage{chngcntr} % fortlaufendes Durchnummerieren der Fußnoten
% \usepackage[perpage]{footmisc} % Alternative: Nummerierung der Fußnoten auf jeder Seite neu

\usepackage{ifthen} % bei der Definition eigener Befehle benötigt
\usepackage{todonotes} % definiert u.a. die Befehle \todo und \listoftodos
\usepackage[square]{natbib} % wichtig für korrekte Zitierweise

% PDF-Optionen -----------------------------------------------------------------
\usepackage{pdfpages}
\pdfminorversion=5 % erlaubt das Einfügen von pdf-Dateien bis Version 1.7, ohne eine Fehlermeldung zu werfen (keine Garantie für fehlerfreies Einbetten!)
\usepackage[
    bookmarks,
    bookmarksnumbered,
    bookmarksopen=true,
    bookmarksopenlevel=1,
    colorlinks=true,
% diese Farbdefinitionen zeichnen Links im PDF farblich aus
    anchorcolor=AOBlau,% Ankertext
    citecolor=AOBlau, % Verweise auf Literaturverzeichniseinträge im Text
    filecolor=AOBlau, % Verknüpfungen, die lokale Dateien öffnen
    menucolor=AOBlau, % Acrobat-Menüpunkte
    urlcolor=AOBlau,
% diese Farbdefinitionen sollten für den Druck verwendet werden (alles schwarz)
    %linkcolor=black, % einfache interne Verknüpfungen
    %anchorcolor=black, % Ankertext
    %citecolor=black, % Verweise auf Literaturverzeichniseinträge im Text
    %filecolor=black, % Verknüpfungen, die lokale Dateien öffnen
    %menucolor=black, % Acrobat-Menüpunkte
    %urlcolor=black,
%
    %backref, % Quellen werden zurück auf ihre Zitate verlinkt
    pdftex,
    plainpages=false, % zur korrekten Erstellung der Bookmarks
    pdfpagelabels=true, % zur korrekten Erstellung der Bookmarks
    hypertexnames=false, % zur korrekten Erstellung der Bookmarks
    linkcolor=black,
    linktoc=all,
]{hyperref}
% Befehle, die Umlaute ausgeben, führen zu Fehlern, wenn sie hyperref als Optionen übergeben werden
\hypersetup{
    pdftitle={\titel -- \untertitel},
    pdfauthor={\autorName},
    pdfcreator={\autorName},
    pdfsubject={\titel -- \untertitel},
    pdfkeywords={\titel -- \untertitel},
}


% zum Einbinden von Programmcode -----------------------------------------------
\usepackage{listings}
\usepackage{xcolor}
\definecolor{hellgelb}{rgb}{1,1,0.9}
\definecolor{colKeys}{rgb}{0,0,1}
\definecolor{colIdentifier}{rgb}{0,0,0}
\definecolor{colComments}{rgb}{0,0.5,0}
\definecolor{colString}{rgb}{1,0,0}
\lstset{
    float=hbp,
	basicstyle=\footnotesize,
    identifierstyle=\color{colIdentifier},
    keywordstyle=\color{colKeys},
    stringstyle=\color{colString},
    commentstyle=\color{colComments},
    backgroundcolor=\color{hellgelb},
    columns=flexible,
    tabsize=2,
    frame=single,
    extendedchars=true,
    showspaces=false,
    showstringspaces=false,
    numbers=left,
    numberstyle=\tiny,
    breaklines=true,
    breakautoindent=true,
	captionpos=b,
}
\lstdefinelanguage{cs}{
	sensitive=false,
	morecomment=[l]{//},
	morecomment=[s]{/*}{*/},
	morestring=[b]",
	morekeywords={
		abstract,event,new,struct,as,explicit,null,switch
		base,extern,object,this,bool,false,operator,throw,
		break,finally,out,true,byte,fixed,override,try,
		case,float,params,typeof,catch,for,private,uint,
		char,foreach,protected,ulong,checked,goto,public,unchecked,
		class,if,readonly,unsafe,const,implicit,ref,ushort,
		continue,in,return,using,decimal,int,sbyte,virtual,
		default,interface,sealed,volatile,delegate,internal,short,void,
		do,is,sizeof,while,double,lock,stackalloc,
		else,long,static,enum,namespace,string},
}
\lstdefinelanguage{natural}{
	sensitive=false,
	morecomment=[l]{/*},
	morestring=[b]",
	morestring=[b]',
	alsodigit={-,*},
	morekeywords={
		DEFINE,DATA,LOCAL,END-DEFINE,WRITE,CALLNAT,PARAMETER,USING,
		IF,NOT,END-IF,ON,*ERROR-NR,ERROR,END-ERROR,ESCAPE,ROUTINE,
		PERFORM,SUBROUTINE,END-SUBROUTINE,CONST,END-FOR,END,FOR,RESIZE,
		ARRAY,TO,BY,VALUE,RESET,COMPRESS,INTO,EQ},
}
\lstdefinelanguage{php}{
	sensitive=false,
	morecomment=[l]{/*},
	morestring=[b]",
	morestring=[b]',
	alsodigit={-,*},
	morekeywords={
		abstract,and,array,as,break,case,catch,cfunction,class,clone,const,
		continue,declare,default,do,else,elseif,enddeclare,endfor,endforeach,
		endif,endswitch,endwhile,extends,final,for,foreach,function,global,
		goto,if,implements,interface,instanceof,namespace,new,old_function,or,
		private,protected,public,static,switch,throw,try,use,var,while,xor
		die,echo,empty,exit,eval,include,include_once,isset,list,require,
		require_once,return,print,unset},
}
 % verwendete Packages
% Sprachdefinition für JavaScript
\lstdefinelanguage{JavaScript}{
  keywords={typeof, new, true, false, catch, function, return, null, catch,
    switch, var, if, in, while, do, else, case, break},
  keywordstyle=\color{blue}\bfseries,
  ndkeywords={class, export, boolean, throw, implements, import, this},
  ndkeywordstyle=\color{darkgray}\bfseries,
  identifierstyle=\color{black},
  sensitive=false,
  comment=[l]{//},
  morecomment=[s]{/*}{*/},
  commentstyle=\color{purple}\ttfamily,
  stringstyle=\color{red}\ttfamily,
  morestring=[b]',
  morestring=[b]"
}
\lstdefinelanguage{AL}
{
  sensitive=false,
  morekeywords=[1]{
    },
  morekeywords=[2]{
    },
  morekeywords=[3]{
    },
  morekeywords=[4]{
    },
  morekeywords=[5]{
    },
  morecomment=[l]{//},
  morecomment=[s]{/*}{*/},
  morestring=[b]"
}

% Stil für AL-Code
\lstset{
    language=AL,
    basicstyle=\small\ttfamily,
    keywordstyle=[1]\color{blue}\bfseries,
    keywordstyle=[2]\color{purple},
    keywordstyle=[3]\color{green!60!black},
    keywordstyle=[4]\color{green!60!black},
    keywordstyle=[5]\color{orange},
    commentstyle=\color{gray},
    stringstyle=\color{orange},
    breaklines=true,
    numbers=left,
    numberstyle=\small,
    frame=tb,
    columns=fullflexible,
    showstringspaces=false,
    tabsize=4,
    backgroundcolor=\color{gray!10},
    captionpos=b
}
% !TEX root = ../Projektdokumentation.tex

% Seitenränder -----------------------------------------------------------------
\setlength{\topskip}{\ht\strutbox} % behebt Warnung von geometry
\geometry{a4paper,left=20mm,right=30mm,top=25mm,bottom=35mm}

\usepackage[
	automark, % Kapitelangaben in Kopfzeile automatisch erstellen
	headsepline, % Trennlinie unter Kopfzeile
	ilines % Trennlinie linksbündig ausrichten
]{scrlayer-scrpage}

% Kopf- und Fußzeilen ----------------------------------------------------------
\pagestyle{scrheadings}
% chapterpagestyle gibt es nicht in scrartcl
%\renewcommand{\chapterpagestyle}{scrheadings}
\clearscrheadfoot

% Kopfzeile
\renewcommand{\headfont}{\normalfont} % Schriftform der Kopfzeile
\ihead{\large{\textsc{\titel}}\\ \small{\untertitel} \\[2ex] \textit{\headmark}}
\chead{}
\setlength{\headheight}{15mm} % Höhe der Kopfzeile
%\setheadwidth[0pt]{textwithmarginpar} % Kopfzeile über den Text hinaus verbreitern (falls Logo den Text überdeckt)

% Fußzeile
\ifoot{\autorName}
\cfoot{}
\ofoot{\pagemark}

% Überschriften nach DIN 5008 in einer Fluchtlinie
% ------------------------------------------------------------------------------

% Abstand zwischen Nummerierung und Überschrift definieren
% > Schön wäre hier die dynamische Berechnung des Abstandes in Abhängigkeit
% > der Verschachtelungstiefe des Inhaltsverzeichnisses
\newcommand{\headingSpace}{1.5cm}

% Abschnittsüberschriften im selben Stil wie beim Inhaltsverzeichnis einrücken
\renewcommand*{\othersectionlevelsformat}[3]{
  \makebox[\headingSpace][l]{#3\autodot}
}

% Für die Einrückung wird das Paket tocloft benötigt
%\cftsetindents{chapter}{0.0cm}{\headingSpace}
\cftsetindents{section}{0.0cm}{\headingSpace}
\cftsetindents{subsection}{0.0cm}{\headingSpace}
\cftsetindents{subsubsection}{0.0cm}{\headingSpace}
\cftsetindents{figure}{0.0cm}{\headingSpace}
\cftsetindents{table}{0.0cm}{\headingSpace}


% Allgemeines
% ------------------------------------------------------------------------------

\onehalfspacing % Zeilenabstand 1,5 Zeilen
\frenchspacing % erzeugt ein wenig mehr Platz hinter einem Punkt

% Schusterjungen und Hurenkinder vermeiden
\clubpenalty = 10000
\widowpenalty = 10000
\displaywidowpenalty = 10000

% Quellcode-Ausgabe formatieren
\lstset{numbers=left, numberstyle=\tiny, numbersep=5pt, breaklines=true}
\lstset{emph={square}, emphstyle=\color{red}, emph={[2]root,base}, emphstyle={[2]\color{blue}}}

\counterwithout{footnote}{section} % Fußnoten fortlaufend durchnummerieren
\setcounter{tocdepth}{3} % im Inhaltsverzeichnis werden die Kapitel bis zum Level der subsubsection übernommen
\setcounter{secnumdepth}{3} % Kapitel bis zum Level der subsubsection werden nummeriert

% Aufzählungen anpassen
\renewcommand{\labelenumi}{\arabic{enumi}.}
\renewcommand{\labelenumii}{\arabic{enumi}.\arabic{enumii}.}
\renewcommand{\labelenumiii}{\arabic{enumi}.\arabic{enumii}.\arabic{enumiii}}

% Tabellenfärbung:
\definecolor{heading}{rgb}{0.64,0.78,0.86}
\definecolor{odd}{rgb}{0.9,0.9,0.9}
 % Definitionen zum Aussehen der Seiten
% !TEX root = ../Projektdokumentation.tex

% Abkürzungen, ggfs. mit korrektem Leerraum
\newcommand{\bs}{$\backslash$\xspace}
\newcommand{\bspw}{bspw.\xspace}
\newcommand{\bzw}{bzw.\xspace}
\newcommand{\ca}{ca.\xspace}
\newcommand{\dahe}{\mbox{d.\,h.}\xspace}
\newcommand{\etc}{etc.\xspace}
\newcommand{\eur}[1]{\mbox{#1\,\texteuro}\xspace}
\newcommand{\evtl}{evtl.\xspace}
\newcommand{\ggfs}{ggfs.\xspace}
\newcommand{\Ggfs}{Ggfs.\xspace}
\newcommand{\gqq}[1]{\glqq{}#1\grqq{}}
\newcommand{\inkl}{inkl.\xspace}
\newcommand{\insb}{insb.\xspace}
\newcommand{\ua}{\mbox{u.\,a.}\xspace}
\newcommand{\usw}{usw.\xspace}
\newcommand{\Vgl}{Vgl.\xspace}
\newcommand{\zB}{\mbox{z.\,B.}\xspace}

% Befehle für häufig anfallende Aufgaben
\newcommand{\Abbildung}[1]{\autoref{fig:#1}}
\newcommand{\Anhang}[1]{\appendixname{}~\ref{#1}: \nameref{#1} \vpageref{#1}}
\newcommand{\includegraphicsKeepAspectRatio}[2]{\includegraphics[width=#2\textwidth,height=#2\textheight,keepaspectratio]{#1}}
\newcommand{\Zitat}[2][\empty]{\ifthenelse{\equal{#1}{\empty}}{\citep{#2}}{\citep[#1]{#2}}}
\newcommand{\Autor}[1]{\textsc{#1}} % zum Ausgeben von Autoren
\newcommand{\itemd}[2]{\item{\textbf{#1}}\\{#2}} % erzeugt ein Listenelement mit fetter Überschrift

% fügt Tabellen aus einer TEX-Datei ein
\newcommand{\tabelle}[3] % Parameter: caption, label, file
{\begin{table}[htbp]
\centering
\singlespacing
\input{Tabellen/#3}
\caption{#1}
\label{#2}
\end{table}}

\newcommand{\tabelleAnhang}[1] % Parameter: file
{\begin{center}
\singlespacing
\input{Tabellen/#1}
\end{center}}

% einfaches Wechseln der Schrift, z.B.: \changefont{cmss}{sbc}{n}
\newcommand{\changefont}[3]{\fontfamily{#1} \fontseries{#2} \fontshape{#3} \selectfont}

% Verwendung analog zu \includegraphics
\newlength{\myx} % Variable zum Speichern der Bildbreite
\newlength{\myy} % Variable zum Speichern der Bildhöhe
\newcommand\includegraphicstotab[2][\relax]{%
% Abspeichern der Bildabmessungen
\settowidth{\myx}{\includegraphics[{#1}]{#2}}%
\settoheight{\myy}{\includegraphics[{#1}]{#2}}%
% das eigentliche Einfügen
\parbox[c][1.1\myy][c]{\myx}{%
\includegraphics[{#1}]{#2}}%
}

\definecolor{AOBlau}{rgb}{0, 0.28, 0.56}

% verschiedene Befehle um Wörter semantisch auszuzeichnen ----------------------
\newcommand{\Index}[2][\empty]{\ifthenelse{\equal{#1}{\empty}}{\index{#2}#2}{\index{#1}#2}}
\newcommand{\Fachbegriff}[2][\empty]{\ifthenelse{\equal{#1}{\empty}}{\textit{\Index{#2}}}{\textit{\Index[#1]{#2}}}}
\newcommand{\NeuerBegriff}[2][\empty]{\ifthenelse{\equal{#1}{\empty}}{\textbf{\Index{#2}}}{\textbf{\Index[#1]{#2}}}}

\newcommand{\Ausgabe}[1]{\texttt{#1}}
\newcommand{\Eingabe}[1]{\texttt{#1}}
\newcommand{\Code}[1]{\texttt{#1}}
\newcommand{\Datei}[1]{\texttt{#1}}

\newcommand{\Assembly}[1]{\textsf{#1}}
\newcommand{\Klasse}[1]{\textsf{#1}}
\newcommand{\Methode}[1]{\textsf{#1}}
\newcommand{\Attribut}[1]{\textsf{#1}}

\newcommand{\Datentyp}[1]{\textsf{#1}}
\newcommand{\XMLElement}[1]{\textsf{#1}}
\newcommand{\Webservice}[1]{\textsf{#1}}

\newcommand{\Refactoring}[1]{\Fachbegriff{#1}}
\newcommand{\CodeSmell}[1]{\Fachbegriff{#1}}
\newcommand{\Metrik}[1]{\Fachbegriff{#1}}
\newcommand{\DesignPattern}[1]{\Fachbegriff{#1}}
 % eigene allgemeine Befehle, die z.B. die Arbeit mit LaTeX erleichtern
% !TEX root = ../Projektdokumentation.tex

% Abkürzungen, ggfs. mit korrektem Leerraum
\newcommand{\bs}{$\backslash$\xspace}
\newcommand{\bspw}{bspw.\xspace}
\newcommand{\bzw}{bzw.\xspace}
\newcommand{\ca}{ca.\xspace}
\newcommand{\dahe}{\mbox{d.\,h.}\xspace}
\newcommand{\etc}{etc.\xspace}
\newcommand{\eur}[1]{\mbox{#1\,\texteuro}\xspace}
\newcommand{\evtl}{evtl.\xspace}
\newcommand{\ggfs}{ggfs.\xspace}
\newcommand{\Ggfs}{Ggfs.\xspace}
\newcommand{\gqq}[1]{\glqq{}#1\grqq{}}
\newcommand{\inkl}{inkl.\xspace}
\newcommand{\insb}{insb.\xspace}
\newcommand{\ua}{\mbox{u.\,a.}\xspace}
\newcommand{\usw}{usw.\xspace}
\newcommand{\Vgl}{Vgl.\xspace}
\newcommand{\zB}{\mbox{z.\,B.}\xspace}

% Befehle für häufig anfallende Aufgaben
\newcommand{\Abbildung}[1]{\autoref{fig:#1}}
\newcommand{\Anhang}[1]{\appendixname{}~\ref{#1}: \nameref{#1} \vpageref{#1}}
\newcommand{\includegraphicsKeepAspectRatio}[2]{\includegraphics[width=#2\textwidth,height=#2\textheight,keepaspectratio]{#1}}
\newcommand{\Zitat}[2][\empty]{\ifthenelse{\equal{#1}{\empty}}{\citep{#2}}{\citep[#1]{#2}}}
\newcommand{\Autor}[1]{\textsc{#1}} % zum Ausgeben von Autoren
\newcommand{\itemd}[2]{\item{\textbf{#1}}\\{#2}} % erzeugt ein Listenelement mit fetter Überschrift

% fügt Tabellen aus einer TEX-Datei ein
\newcommand{\tabelle}[3] % Parameter: caption, label, file
{\begin{table}[htbp]
\centering
\singlespacing
\input{Tabellen/#3}
\caption{#1}
\label{#2}
\end{table}}

\newcommand{\tabelleAnhang}[1] % Parameter: file
{\begin{center}
\singlespacing
\input{Tabellen/#1}
\end{center}}

% einfaches Wechseln der Schrift, z.B.: \changefont{cmss}{sbc}{n}
\newcommand{\changefont}[3]{\fontfamily{#1} \fontseries{#2} \fontshape{#3} \selectfont}

% Verwendung analog zu \includegraphics
\newlength{\myx} % Variable zum Speichern der Bildbreite
\newlength{\myy} % Variable zum Speichern der Bildhöhe
\newcommand\includegraphicstotab[2][\relax]{%
% Abspeichern der Bildabmessungen
\settowidth{\myx}{\includegraphics[{#1}]{#2}}%
\settoheight{\myy}{\includegraphics[{#1}]{#2}}%
% das eigentliche Einfügen
\parbox[c][1.1\myy][c]{\myx}{%
\includegraphics[{#1}]{#2}}%
}

\definecolor{AOBlau}{rgb}{0, 0.28, 0.56}

% verschiedene Befehle um Wörter semantisch auszuzeichnen ----------------------
\newcommand{\Index}[2][\empty]{\ifthenelse{\equal{#1}{\empty}}{\index{#2}#2}{\index{#1}#2}}
\newcommand{\Fachbegriff}[2][\empty]{\ifthenelse{\equal{#1}{\empty}}{\textit{\Index{#2}}}{\textit{\Index[#1]{#2}}}}
\newcommand{\NeuerBegriff}[2][\empty]{\ifthenelse{\equal{#1}{\empty}}{\textbf{\Index{#2}}}{\textbf{\Index[#1]{#2}}}}

\newcommand{\Ausgabe}[1]{\texttt{#1}}
\newcommand{\Eingabe}[1]{\texttt{#1}}
\newcommand{\Code}[1]{\texttt{#1}}
\newcommand{\Datei}[1]{\texttt{#1}}

\newcommand{\Assembly}[1]{\textsf{#1}}
\newcommand{\Klasse}[1]{\textsf{#1}}
\newcommand{\Methode}[1]{\textsf{#1}}
\newcommand{\Attribut}[1]{\textsf{#1}}

\newcommand{\Datentyp}[1]{\textsf{#1}}
\newcommand{\XMLElement}[1]{\textsf{#1}}
\newcommand{\Webservice}[1]{\textsf{#1}}

\newcommand{\Refactoring}[1]{\Fachbegriff{#1}}
\newcommand{\CodeSmell}[1]{\Fachbegriff{#1}}
\newcommand{\Metrik}[1]{\Fachbegriff{#1}}
\newcommand{\DesignPattern}[1]{\Fachbegriff{#1}}
 % eigene projektspezifische Befehle, z.B. Abkürzungen usw.

\begin{document}
% ---------------------------------------------------------------------------
%\phantomsection
%\thispagestyle{plain}
%\pdfbookmark[1]{Eidesstattliche Erklärung}{ihkdeckblatt}
%\includegraphicsKeepAspectRatio{DeckblattIHK.pdf}{1}
%\cleardoublepage

\phantomsection
\thispagestyle{plain}
\pdfbookmark[1]{Deckblatt}{deckblatt}
% !TEX root = Projektdokumentation.tex
\begin{titlepage}

\begin{center}
\includegraphics[scale=0.25]{LogoIHK.pdf}\\[1ex]
\Large{Abschlussprüfung \pruefungstermin}\\[3ex]

\Large{\ausbildungsberuf}\\
\LARGE{\betreff}\\[4ex]

\huge{\textbf{\titel}}\\[1.5ex]
\Large{\textbf{\untertitel}}\\[4ex]

\normalsize
Abgabetermin: \abgabeOrt, den \abgabeTermin\\[3em]
\textbf{Prüfungsbewerber:}\\
\autorName\\
\autorAnschrift\\
\autorOrt\\[5ex]

\includegraphics[scale=0.4]{\betriebLogo}\\[2ex]
\textbf{Ausbildungsbetrieb:}\\
\betriebName\\
\betriebAnschrift\\
\betriebOrt\\[5em]
\end{center}

\end{titlepage}
\cleardoublepage

% Preface --------------------------------------------------------------------
\phantomsection
\pagenumbering{Roman}
\pdfbookmark[1]{Inhaltsverzeichnis}{inhalt}
\tableofcontents

\cleardoublepage

\phantomsection
\listoffigures
\cleardoublepage

\phantomsection
\listoftables
\cleardoublepage

%\phantomsection
%\lstlistoflistings
%\cleardoublepage

\newcommand{\abkvz}{Abkürzungsverzeichnis}
\renewcommand{\nomname}{\abkvz}
\section*{\abkvz}
\markboth{\abkvz}{\abkvz}
\addcontentsline{toc}{section}{\abkvz}
% !TEX root = Projektdokumentation.tex

% Es werden nur die Abkürzungen aufgelistet, die mit \ac definiert und auch benutzt wurden. 
%
% \acro{VERSIS}{Versicherungsinformationssystem\acroextra{ (Bestandsführungssystem)}}
% Ergibt in der Liste: VERSIS Versicherungsinformationssystem (Bestandsführungssystem)
% Im Text aber: \ac{VERSIS} -> Versicherungsinformationssystem (VERSIS)

% Hinweis: allgemein bekannte Abkürzungen wie z.B. bzw. u.a. müssen nicht ins Abkürzungsverzeichnis aufgenommen werden
% Hinweis: allgemein bekannte IT-Begriffe wie Datenbank oder Programmiersprache müssen nicht erläutert werden,
%          aber ggfs. Fachbegriffe aus der Domäne des Prüflings (z.B. Versicherung)

% Die Option (in den eckigen Klammern) enthält das längste Label oder
% einen Platzhalter der die Breite der linken Spalte bestimmt.
\begin{acronym}[WWWWW]
    \acro{AL-Code}{Application Language Code}
    \acro{BC}{Business Central, Software für ERP-Systeme, welches für Warenwirtschaft, Vertrieb, Service und Finanzen genutzt wird}
    \acro{CRM}{Customer Relationship Management}
    \acro{Dynamics}{Microsoft Dynamics CRM ist eine leistungsstarke CRM-Plattform, die Sie dabei unterstützt, den Vertriebsprozess zu automatisieren und zu optimieren}
    \acro{ERP}{Enterprise Resource Planning}
    \acro{NAS}{Network Attached Storage, einfach zu verwaltender Dateiserver}
    \acro{SCM}{Supply Chain Management}
    \acro{API}{Application Programming Interface}   
    \acro{CSV}{Comma Separated Values}
    \acro{JSON}{JavaScript Object Notation}
    \acro{VM}{VONMÄHLEN}
    \acro{SQL}{Structured Query Language}
    \acro{ZF}{Zeitfracht}
    \acro{MVC}{Model-View-Controller}   
    \acro{IDE}{Integrated Development Environment}
    \acro{Dbeaver}{Datenbankverwaltungstool}  
\end{acronym}




\clearpage

% Inhalt ---------------------------------------------------------------------
\pagenumbering{arabic}
% !TEX root = Projektdokumentation.tex
% !TEX root = ../Projektdokumentation.tex
\section{Einleitung}
\label{sec:Einleitung}

\subsection{Projektumfeld} 
\label{sec:Projektumfeld}
VONMÄHLEN ist ein deutsches Unternehmen, das sich auf die Entwicklung und Produktion von hochwertigen, 
designorientierten Lifestyle-Technologieprodukten spezialisiert hat. Die Firma bietet innovative Zubehörlösungen für den täglichen Gebrauch, 
insbesondere im Bereich Smartphone und Technik-Accessoires.

\subsection{Projektziel} 
\label{sec:Projektziel}
Die SCM-Abteilung wünscht sich eine automatisierte Überwachung der Warensendungen, um den manuellen Aufwand und das Risiko menschlicher Fehler zu minimieren. Das System soll den Status aller Warensendungen automatisch und regelmäßig überprüfen und bei Abweichungen oder Verzögerungen eine E-Mail-Benachrichtigung an die zuständigen Mitarbeiter der SCM-Abteilung senden.
Zusätzlich soll das System im ERP-System (Microsoft Business Central) den Zeitpunkt der Auslieferung einer Sendung erfassen und aktualisieren, damit ersichtlich ist, wann eine Lieferung erfolgreich zugestellt wurde. Diese Information dient als Grundlage dafür, dass die SCM-Abteilung die Buchung der zugehörigen Rechnung anstoßen kann, sobald die Lieferung den Empfänger erreicht hat. 
Durch diese Automatisierung erwartet die SCM-Abteilung eine höhere Effizienz, da die Zustellvorgänge transparent und leicht nachvollziehbar sind und die Rechnungsstellung unmittelbar nach der Zustellung erfolgen kann.
Eine detaillierte Auflistung der Anforderungen ist dem Lastenheft zu entnehmen. 
Siehe \Anhang{app:Lastenheft}


\subsection{Projektbegründung} 
\label{sec:Projektbegruendung}
Die größte Schwachstelle im aktuellen Prozess liegt im hohen manuellen Aufwand und der fehlenden Dokumentation. 
Täglich werden alle offenen Lieferaufträge nach solchen mit einer Trackingnummer durchsucht, was als Indikator dient, 
dass die Ware bereits auf dem Weg ist und überprüft werden muss. Dabei wird festgestellt, wie weit die Ware ist und ob 
Verzögerungen oder Probleme vorliegen. Die Prüfung erfolgt durch Eingabe der Trackingnummer im Webportal des 
Versanddienstleisters (z. B. DPD). Je nach Lieferstatus wird dann entsprechend gehandelt, und bei Auffälligkeiten 
wird die zuständige Abteilung informiert. Dieser manuelle Prozess birgt jedoch das Risiko, dass Überprüfungen 
ausgelassen oder ganze Aufträge übersehen werden. Auch kann es bei Personalmangel zum vollständigen Ausfall der 
Prüfungen kommen. Aufgrund dieser Fehleranfälligkeit und des hohen Zeitaufwands hat sich die Firma VONMÄHLEN 
entschieden, den Prozess zu automatisieren, um die Effizienz zu steigern und die Prozesssicherheit nachhaltig zu 
gewährleisten.

\subsection{Technische Projektschnittstellen} 
\label{sec:Projektschnittstellen}
Diese Schnittstellen gewährleisten eine reibungslose Interaktion zwischen den verschiedenen Systemen und Komponenten, 
die für das Projekt erforderlich sind.

\subsubsection{Versanddienstleister: Zeitfracht (ZF)}
\begin{itemize}
\item Art: SFTP-Schnittstelle
\item Zweck: Bereitstellung von Trackingdaten
\item Beschreibung: Der Versanddienstleister liefert täglich aktualisierte Tracking-Daten als CSV-Dateien über einen 
gesicherten SFTP-Server. Diese Daten werden vom System heruntergeladen, validiert und verarbeitet.
\end{itemize}

\subsubsection{ERP-System: Microsoft Business Central}
\begin{itemize}
\item Art: API-Schnittstelle
\item Zweck: Austausch von Informationen zu Sendungsstatus und Aktualisierungen von Lieferdaten
\item Beschreibung: Die API ermöglicht es, Statusinformationen und Lieferdetails aus der 
PostgreSQL-Datenbank in das ERP-System zu übertragen, um die automatisierte Weiterverarbeitung 
(z.B. Rechnungsstellung) zu gewährleisten.
\end{itemize}

\subsubsection{Datenbank: PostgreSQL}
\begin{itemize}
\item Art: SQL-Schnittstelle
\item Zweck: Speicherung und Abfrage der Tracking-Daten
\item Beschreibung: Diese Schnittstellen gewährleisten eine reibungslose Interaktion zwischen den verschiedenen Systemen und Komponenten, die für das Projekt erforderlich sind.
\end{itemize}

\subsubsection{Workflow-Management-System: Apache Airflow}
\begin{itemize}
\item Art: Workflow-Orchestrierungs-Tool
\item Zweck: Automatisierung und Orchestrierung von Prozessen
\item Beschreibung: Apache Airflow steuert die zeitgesteuerte Verarbeitung der Trackingdaten, die Kommunikation mit der Datenbank und die Synchronisierung mit dem ERP-System.
\end{itemize}

\subsubsection{Docker-Container}
\begin{itemize}
\item Art: Container-Übergreifende Kommunikation
\item Zweck: Isolierte Ausführung von PostgreSQL und Apache Airflow
\item Beschreibung: Docker-Container werden genutzt, um PostgreSQL und Apache Airflow in einer sicheren und skalierbaren Umgebung auszuführen. Die Container sind so konfiguriert, dass sie miteinander kommunizieren können, um Daten auszutauschen. 
\end{itemize}

% !TEX root = ../Projektdokumentation.tex
\section{Projektplanung} 

\subsection{Projektphasen}
\label{sec:Projektphasen}
Für diese Umsetzung des Projektes standen dem Auszubildenen 80 Stunden zur Verfügung. 
Welche vor Projektbeginn in verschiedene Phasen aufgeteilt wurden. 

Tabelle~\ref{tab:Zeitplanung} zeigt eine grobe Zeitplanung.
\tabelle{Zeitplanung}{tab:Zeitplanung}{ZeitplanungKurz}\\

Eine detailliertere Zeitplanung ist im \Anhang{app:Zeitplanung}

\clearpage


\subsection{Abweichungen vom Projektantrag}
\label{sec:Abweichung Projektantrag}
Es stellte sich während des Projektes heraus das einige Zeiten nicht eingehalten werden konnten. 
Darauf wird im \Anhang{app:Zeitnachher} weiter eingegangen.


\subsection{Ressourcenplanung}
\label{sec:Ressourcenplanung}
Um die Planung der erforderlichen Ressourcen effizient zu gestalten, wurde gezielt auf bereits vorhandene Infrastruktur zurückgegriffen. Die Auswahl der Hardware und Software erfolgte so, dass keine zusätzlichen Anschaffungen erforderlich waren. Hardware war vollständig vorhanden, da alle Mitarbeiter über einen ausgestatteten Büroarbeitsplatz mit Standard-Peripherie verfügten. Die Softwareauswahl konzentrierte sich auf Open-Source-Tools wie PostgreSQL und Python sowie bereits lizenzierte Programme (Business Central API), die während der Projektlaufzeit genutzt wurden, um die Datenverarbeitung sicherzustellen.
Der Auszubildende übernahm die Hauptverantwortung in der Entwicklung, was insgesamt 80 Stunden veranschlagte. Ein Entwickler wurde für 10 Stunden eingeplant, um gezielt bei technischen Fragen zu unterstützen. Für die infrastrukturelle Vorbereitung, wie das Aufsetzen der Docker-Container inklusive der PostgreSQL-Datenbank und Apache Airflow, wurde Zuarbeit durch die IT-Abteilung geleistet. Dies stellte sicher, dass eine stabile und funktionale Grundlage für die Entwicklung geschaffen wurde.
Die SCM-Abteilung war schließlich dafür verantwortlich, die Anforderungen zu formulieren und das System während der Testphase auf seine Praxistauglichkeit zu überprüfen. Die Kombination aus vorhandenen Ressourcen und gezielter, bedarfsorientierter Planung sorgte dafür, dass die Projektkosten gering blieben und die notwendige Unterstützung jederzeit gewährleistet war.


\subsection{Entwicklungsprozess}
\label{sec:Entwicklungsprozess}
Für die Entwicklung dieses Projekts wurde das Wasserfallmodell als Vorgehensmodell gewählt. 
Das Wasserfallmodell zeichnet sich durch eine klare und lineare Abfolge von Phasen aus, was insbesondere bei Projekten mit fest 
definierten Anforderungen und einem strukturierten Ablauf von Vorteil ist, zudem ist es ein Vorgehensmodell der klassischen Softwareentwicklung, bei dem die Projektphasen sequentiell 
durchlaufen und bearbeitet werden (vgl. \cite{FittkauRuf2008} S. 31). 
Da die Anforderungen an das System im Vorfeld detailliert analysiert und dokumentiert wurden, ermöglichte das Wasserfallmodell 
eine schrittweise und systematische Umsetzung der Projektphasen von der Anforderungsanalyse über die Implementierung bis hin zur abschließenden Testphase.


% !TEX root = ../Projektdokumentation.tex
\section{Analysephase} 
\label{sec:Analysephase}

\subsection{Ist-Analyse} 
\label{sec:IstAnalyse}
Derzeit werden täglich alle Warensendungen manuell geprüft, indem Lieferungen mit Trackingnummern 
herausgefiltert und im Onlineportal des externen Versanddienstleisters überprüft werden. 
Mitarbeiter in der Abteilung \ac{SCM} durchforsten Verkaufsaufträge nach 
Trackingnummern und prüfen im Versandportal den aktuellen Status der Sendungen. Der Fokus liegt 
hierbei darauf, Unregelmäßigkeiten wie Verzögerungen oder Fehlzustellungen zu identifizieren. 
Bei auffälligen Zuständen wird die zuständige Abteilung informiert, um die erforderlichen Maßnahmen einzuleiten.

\subsection{Wirtschaftlichkeitsanalyse}
\label{sec:Wirtschaftlichkeitsanalyse}
Die Wirtschaftlichkeit des Projekts ließ sich unkompliziert anhand der eingesparten Arbeitszeit der SCM-Mitarbeiter berechnen, 
da diese durch die Automatisierung des Prüfungsprozesses deutlich reduziert wird.


\subsection{Kostenanalyse}
\label{sec:Kostenanalyse}
Die Kosten für die Implementierung des Systems umfassen zwei Hauptkategorien: 
Initiale Implementierungskosten und laufende Betriebskosten.

\subsubsection{Einmalige Implementierungskosten}
\label{sec:Implementierungskosten}

Diese Kosten fallen nur einmal an und umfassen die Entwicklung des Systems, 
die Anschaffung notwendiger Software sowie die Einrichtung der Infrastruktur.

	\textbf{Entwicklungskosten (inklusive Planung, Design und Implementierung)}:
    \begin{itemize}
        \item Auszubildender: 80 Stunden a 7 €/Stunde = 560 €
        \item Entwickler-Support: 10 Stunden a 22 €/Stunde = 220 €
        \item Fachabteilung SCM: 2 Stunden a 30 €/Stunde
    \end{itemize}
    \clearpage

	\textbf{Softwarelizenzen und Tools:}
    \begin{itemize}
        \item PostgreSQL: Kostenlos
        \item Visual Studio Code: Kostenlos
        \item Docker: Kostenlos
        \item Filezilla: Kostenlos
        \item Apache Airflow: Kostenlos
        \item Business Central Lizenz (während der Implementierung): 65,50 €/Monat, über 2 Monate = 131 € 
    \end{itemize}

	\textbf{Hardware- und Infrastrukturkosten:}
    \begin{itemize}
        \item Serverkosten, Server vorhanden: 0€/Jahr
    \end{itemize}
        
	\textbf{Gesamtkosten für die Implementierung: 971 €}

\subsubsection{Laufende Betriebskosten}
\label{Betriebskosten}

Diese Kosten fallen während des Betriebs des Systems an.

\textbf{Wartung, Fehlerbehebung und Optimierungen:}
\begin{itemize}
    \item Entwickler-Support (12 Stunden pro Jahr a 22€/Stunde) = 264 €/Jahr
    \item Erweiterung des bestehenden Hetzner Servers 120 €/Jahr
\end{itemize}

\textbf{Gesamtkosten für den laufenden Betrieb: 384 €/Jahr }

\clearpage

\subsection{Vergleich: Manuelle vs. Automatisierte Prozesse}
\label{vergleichProzesse}
Tabelle~\ref{tab:vergleichProzesse}
\tabelle{Vergleich von manuellen und automatisierten Prozessen}{tab:vergleichProzesse}{vergleichProzesse}\\




















% !TEX root = ../Projektdokumentation.tex
\section{Entwurfsphase} 
\label{sec:Entwurfsphase}

\subsection{Zielplattform}
\label{sec:Zielplattform} 
Die Zielplattform für dieses Projekt ist die IT-Infrastruktur der VONMÄHLEN GmbH, die auf ein internes Netzwerk 
ausgerichtet ist. Kernkomponenten des Projekts, darunter die PostgreSQL-Datenbank und das \ac{ERP}-System BC, sind in 
die bestehende Serverarchitektur integriert. Zur Automatisierung von Aufgaben wie dem täglichen Datenabgleich und 
der Aktualisierung der Trackinginformationen werden Docker und Apache Airflow eingesetzt. Diese Werkzeuge ermöglichen 
eine nahtlose Integration und die zuverlässige Ausführung wieder-kehrender Aufgaben. Der Zugriff auf Business Central 
erfolgt über eine abgesicherte \ac{API}-Schnittstelle, während der Datentransfer über einen SFTP-Server abgewickelt wird, 
der die Trackingdaten vom externen Dienstleister sicher aufnimmt. Mit dieser Infrastruktur wird eine optimale Sicherheit 
und Performance erreicht, und gleichzeitig können die Anforderungen an eine hohe Verfügbarkeit und Skalierbarkeit erfüllt 
werden. Architekturdesign Beschreibung und Begründung der gewählten Anwendungsarchitektur \ac{MVC}. Ggfs. Bewertung und 
Auswahl von verwendeten Frameworks sowie ggfs. eine kurze Einführung in die Funktionsweise des verwendeten Frameworks.

\subsection{Architekturdesign}
\label{sec:Architekturdesign}
Python in Verbindung mit Apache Airflow wurde aufgrund der einfachen Integration und Technologiengleichheit im Unternehmen gewählt. 
Python ermöglicht schnelle Entwicklung, während Airflow die Orchestrierung und Überwachung der Datenpipelines übernimmt. Ideal für 
die automatisierte Verarbeitung der Tracking-Daten.


% !TEX root = ../Projektdokumentation.tex
\section{Implementierungsphase} 
\label{sec:Implementierungsphase}

\subsection{PostgreSQL-Datenbankmodellierung und -Erstellung}
\label{sec:Datenbankmodellierung}

Die PostgreSQL-Datenbank umfasst eine einzige Tabelle namens \textit{bc\_tracking}. 
Aufgrund der einfachen Projektanforderungen war ein komplexes ER-Modell nicht nötig. 
Die Datenbank speichert Tracking-Informationen von Zeitfracht und BC.

Tabelle~\ref{tab:bcTracking} Die Struktur der Tabelle \textit{bc\_tracking} ist wie folgt:
\tabelle{\textit{bc\_tracking}}{tab:bcTracking}{bcTracking}\\

\begin{figure}[htb]
    \centering
    \includegraphicsKeepAspectRatio{sql_create_db.png}{0.5}
    \caption{Skript zur Erstellung der Datenbank}
    \label{fig:DatenbankBCTracking}
\end{figure}

\clearpage

\subsection{5.2	Deployment auf dem Hetzner-Server}
\label{sec:Hetzner-Server}

Die PostgreSQL-Datenbank wurde auf einem Hetzner-Server installiert, um eine zuverlässige und performante 
Umgebung für die Anwendung bereitzustellen. Die Verwaltung der Datenbank erfolgt über DBeaver, das sowohl 
lokale als auch entfernte Datenbankverbindungen unterstützt.

\subsection{Testdaten und Verifizierung des Datenbankzugriffs}
\label{sec:VerifizierungTestdaten}

Um die korrekte Funktion der Datenbank und der Anwendungen zu verifizieren, wurden Test-daten verwendet. 
Diese Testdaten wurden in die Tabelle \textit{bc\_tracking} eingefügt und die An-wendungen darauf ausgeführt, 
um sicherzustellen, dass die Daten korrekt verarbeitet und die gewünschten Aktionen, wie das Senden von E-Mails, 
ausgeführt werden. Screenshots der Anwendung in der Entwicklungsphase mit Dummy-Daten befinden sich im Anhang A3 Abbildung 7

\subsection{Implementierung der Geschäftslogik}
\label{sec:Geschaeftslogik}

Die Geschäftslogik wurde in Python implementiert und besteht aus zwei Hauptskripten: trackings.py und \textit{bc\_db.py}. 
Das Skript trackings.py ist verantwortlich für das Einlesen von CSV-Dateien mit Tracking-Daten von Zeitfracht und das 
Speichern dieser Daten in der PostgreSQL-Datenbank. Das Skript führt folgende Schritte aus:
\begin{itemize}
    \item Einlesen aller CSV-Dateien aus einem Quellverzeichnis.
    \item Validierung der Dateiinhalte, insbesondere der notwendigen Spalten.
    \item Verarbeitung jeder Zeile und Einfügen der Daten in die Tabelle \textit{bc\_tracking}.
    \item Bei erfolgreicher Verarbeitung wird die CSV-Datei aus dem Quellverzeichnis gelöscht.
\end{itemize}


Ein wichtiger Bestandteil ist die Validierung der Daten. Hier wird sichergestellt, dass nur gültige
Statuscodes verarbeitet werden:	

\begin{figure}[htb]
    \centering
    \includegraphicsKeepAspectRatio{validate_row.png}{1}
    \caption{Validierung der Statuscodes}
    \label{fig:validateRow}
\end{figure}

Das Skript \textit{bc\_db.py} kümmert sich um die Verarbeitung der in der Datenbank gespeicherten Trackingdaten und 
interagiert mit der Business Central API. 
Die Hauptaufgaben sind:

\begin{itemize}
    \item Überprüfen, ob bestimmte Tracking-Nummern länger als eine definierte Anzahl von Tagen in einem bestimmten Status verweilen.
    \item Senden von Benachrichtigungs-E-Mails, falls diese Bedingungen erfüllt sind
    \item Senden von Benachrichtigungs-E-Mails, falls Fehler auftreten
    \item Aktualisieren von Datensätzen in Business Central über die API.
\end{itemize}


Beispielsweise wird geprüft, ob Sendungen mit Statuscode 30 länger als 7 Tage nicht zugestellt wurden:

\begin{figure}[htb]
    \centering
    \includegraphicsKeepAspectRatio{sql_view.png}{1}
    \caption{Prüfung von Sendungen im Status 30}
    \label{fig:viewStatus}
\end{figure}


Sowohl \textit{trackings.py} als auch \textit{bc\_db.py} nutzen verschiedene Integrationen, um die Funktionalität zu erweitern. 
Zum Beispiel wird das Decorator-Pattern verwendet, um Integrationen einzubinden:

\begin{figure}[htb]
    \centering
    \includegraphicsKeepAspectRatio{dekorator_pattern.png}{0.8}
    \caption{Verwendung von Integrationen}
    \label{fig:Integrationen}
\end{figure}


Dieses Muster ermöglicht eine modulare und erweiterbare Architektur der Anwendung.

\clearpage

Des Weiteren werden Ausnahmen und Fehler sorgfältig behandelt, um die Stabilität der An-wendung zu gewährleisten. 
Zum Beispiel wird beim Aktualisieren von Daten in BC auf mögliche API-Fehler reagiert:

\begin{figure}[htb]
    \centering
    \includegraphicsKeepAspectRatio{update_data.png}{0.8}
    \caption{Fehlerbehandlung bei API-Aufrufen}
    \label{fig:viewAPI}
\end{figure}

Die Klasse \textit{Zf\_bc\_tracking} in \textit{rackings.py} implementiert die Logik zum Einlesen und Verarbeiten der Trackingdaten 
von Zeitfracht. Zusammenfassend wurde die Geschäftslogik unter Verwendung von Python und verschiedenen Bibliotheken wie SQLAlchemy 
für den Datenbankzugriff, Logging für die Protokollierung und benutzerdefinierten Integrationen für die Anbindung externer Systeme realisiert.



























% !TEX root = ../Projektdokumentation.tex
\section{Abnahmephase}
\label{sec:Abnahmephase}

\subsection{Testphase}
\label{sec:Testphase}
Alle Systeme wurden unter realistischen Bedingungen in einer speziell dafür eingerichteten Entwicklungsumgebung getestet.
Diese Umgebung, die regelmäßig für Integration-Tests genutzt wird, stellt sicher, dass die Systeme stabil und
zuverlässig funktionieren.

Zum Testsystem gehören unter anderem:

\begin{itemize}
	\item Microsoft Dynamics 365 Business Central
	\item PostgreSQL-Datenbank (siehe Anhang)
	\item Apache Airflow für automatisierte Python-Skripte (siehe Anhang)
	\item SFTP-Server für den Datentransfer
\end{itemize}

\subsubsection{Tests mit dem Skript \textit{trackings.py}}
\label{sec:Tests mit dem Skript trackings.py}

Die Tests bezogen sich hauptsächlich auf die folgenden Aspekte:

\begin{enumerate}
    \item Upload von Dummy-Daten auf den SFTP-Server
    \begin{itemize}
        \item Im Rahmen der Tests wurden regelmäßig Dummy-Daten auf den SFTP-Server geladen, um den vollständigen Prozess von der Datenerfassung bis zur Weiterverarbeitung zu simulieren.
        \item Diese Dummy-Daten bestanden aus korrekt formatierten Tracking-Datensätzen sowie absichtlich fehlerhaften Datensätzen, um die Robustheit der Validierungslogik zu prüfen.
    \end{itemize}
    \item Abholung und Validierung der Daten durch das Skript
    \begin{itemize}
        \item Das Skript \textit{trackings.py} wurde verwendet, um die hochgeladenen Dummy-Daten vom SFTP-Server abzuholen und in die PostgreSQL-Datenbank zu speichern.
        \item Eine umfassende Validierung der Daten wurde durchgeführt, um sicherzustellen, dass nur vollständige und korrekte Datensätze verarbeitet werden.
        \item Es wurden insbesondere folgende Validierungsszenarien getestet:
        \begin{itemize}
            \item Vollständige und korrekte Datensätze: Prüfung, ob die Daten korrekt in die Datenbank eingefügt werden.
            \item Fehlerhafte oder unvollständige Datensätze: Sicherstellung, dass fehlerhafte Datensätze (z.B. fehlende Pflichtfelder, falsche Datentypen) nicht in die Datenbank übernommen werden, sondern eine entsprechende Fehlermeldung generieren.
            \item Doppelte Datensätze: Dummy-Daten, die doppelt vorhanden waren, wurden hochgeladen, um sicherzustellen, dass das Skript diese erkennt und nicht erneut in die Datenbank einfügt.
        \end{itemize}
    \end{itemize}
    \item Löschen der erfolgreich gelesenen Daten auf dem SFTP-Server
    \begin{itemize}
        \item Nach der erfolgreichen Verarbeitung der Dummy-Daten wurde das automatische Löschen der verarbeiteten Dateien auf dem SFTP-Server getestet.
        \item Es wurde sichergestellt, dass nur die Daten, die erfolgreich verarbeitet wurden, gelöscht wurden, während nicht verarbeitete oder fehlerhafte Dateien weiterhin auf dem SFTP-Server verbleiben.
        \item Das Skript wurde darauf getestet, korrekt zwischen verarbeiteten und nicht verarbeiteten Dateien zu unterscheiden und die Löschvorgänge nur bei Erfolg auszuführen.
    \end{itemize}
    \item Fehlerhafte Daten gezielt einfügen
    \begin{itemize}
        \item Es wurden gezielt falsche Daten auf den SFTP-Server hochgeladen, um zu prüfen, wie die Skripte auf inkonsistente Informationen reagieren.
        \item Ziel dieser Tests war es, sicherzustellen, dass das System solche fehlerhaften Daten erkennt und entsprechend darauf reagiert, z.B. durch Protokollierung im Log oder durch Benachrichtigung der zuständigen Personen.
    \end{itemize}
    \item Verifizierung der Datenintegrität
    \begin{itemize}
        \item Die Testdaten wurden genutzt, um die Datenintegrität in der PostgreSQL-Datenbank zu überprüfen.
        \item Es wurde sichergestellt, dass nur die validierten Daten korrekt in der Datenbank gespeichert wurden und dass keine fehlerhaften Daten Einfluss auf die vorhandenen Datensätze hatten.
        \item Dies beinhaltete auch die Konsistenzprüfung der Datenbank nach der Verarbeitung.
    \end{itemize}
\end{enumerate}

\subsubsection{Tests mit dem Skript \textit{bc\_db.py}}
\label{sec:Tests mit dem Skript bc_db.py}

Die Testphase beinhaltete die folgenden Aspekte:

\begin{enumerate}
    \item Abfrage und Verarbeitung der Tracking-Daten
    \begin{itemize}
        \item Das Skript \textit{bc\_db.py} wurde getestet, um sicherzustellen, dass es in der Lage ist, die Tracking-Daten korrekt aus der PostgreSQL-Datenbank zu lesen und zu verarbeiten.
        \item Es wurden mehrere Testläufe durchgeführt, um sicherzustellen, dass das Skript die relevanten Daten korrekt abfragt und für die Weiterverarbeitung vorbereitet.
        \item Verarbeitungsszenarien: Verschiedene Szenarien, wie z.B. Datensätze mit unterschiedlichen Statuscodes, wurden geprüft, um sicherzustellen, dass das System die jeweilige Geschäftslogik korrekt anwendet.
    \end{itemize}
    \item Synchronisation mit Microsoft Business Central
    \begin{itemize}
        \item Ein wesentlicher Bestandteil des Skripts ist die Synchronisation der Tracking-Daten mit dem ERP-System Microsoft Business Central.
        \item Die Funktionalität wurde getestet, indem die API-Aufrufe mit verschiedenen Testdaten ausgeführt wurden.
        \item Es wurde insbesondere geprüft, dass die Statusinformationen korrekt in Business Central aktualisiert werden und dass die API fehlerfrei funktioniert.
        \item Hierzu wurden auch absichtlich falsche API-Daten verwendet, um sicherzustellen, dass das System korrekt auf Fehlermeldungen reagiert und diese verarbeitet.
    \end{itemize}
    \item Fehlerbehandlung und Benachrichtigung
    \begin{itemize}
        \item Das Skript wurde getestet, um sicherzustellen, dass es robust gegenüber Fehlern ist, die bei der Verarbeitung der Daten auftreten können.
        \item Es wurden Szenarien simuliert, in denen die API-Verbindung fehlschlägt, um zu prüfen, wie das Skript darauf reagiert.
        \item Zudem wurde die Benachrichtigungslogik getestet: Wenn eine Sendung länger als die definierte Anzahl an Tagen in einem bestimmten Status verweilt, wird automatisch eine Benachrichtigung an die zuständigen Mitarbeiter gesendet.
        \item Die Testläufe stellten sicher, dass diese E-Mails korrekt ausgelöst werden.
    \end{itemize}
    \item Korrekte Handhabung von Ausnahmefällen
    \begin{itemize}
        \item Es wurden Tests durchgeführt, um sicherzustellen, dass alle möglichen Ausnahmefälle korrekt behandelt werden, z.B. wenn die Datenbankverbindung vorübergehend nicht verfügbar ist oder unerwartete Datenformate auf der API-Ebene auftreten.
        \item Die Protokollierung der Fehler wurde ebenfalls geprüft, um sicherzustellen, dass alle Fehler im Log korrekt vermerkt werden, damit diese bei Bedarf nachvollzogen und behoben werden können.
    \end{itemize}
    \item End-to-End-Test der gesamten Datenpipeline
    \begin{itemize}
        \item Schließlich wurde ein End-to-End-Test der gesamten Datenpipeline durchgeführt, bei dem beide Skripte (\textit{trackings.py} und \textit{bc\_db.py}) zusammenarbeiteten.
        \item Ziel war es sicherzustellen, dass die Daten von der Erfassung über den SFTP-Server, die Validierung und Speicherung in der PostgreSQL-Datenbank bis hin zur Synchronisation mit Business Central nahtlos verarbeitet werden.
        \item Diese Tests beinhalteten auch die Simulation realer Arbeitsbedingungen, um die Performance und Zuverlässigkeit des Systems unter hoher Last zu prüfen.
    \end{itemize}
\end{enumerate}

\input{Inhalt/Einführungsphase}
% !TEX root = ../Projektdokumentation.tex
\section{Dokumentation}
\label{sec:Dokumentation}


\subsection{Zielgruppe der Entwicklerdokumentation}
\label{sec:ZielgruppeDoku}

Die Dokumentation richtet sich an administratives Personal der VONMÄHLEN GmbH. 
Der Umgang mit Apache Airflow und PostgreSQL ist diesem Personenkreis geläufig. 
Python wird als unternehmensübliche Programmiersprache genutzt. Die Dokumentation ist 
sowohl im Markdown-Format im GIT-Repository (\textit{README.md}) als auch im internen Unternehmens-Wiki verfügbar.

\subsection{Inhalt der Entwicklerdokumentation}
\label{sec:Entwicklerdokumentation}

Die Entwicklerdokumentation umfasst:
\begin{itemize}
    \item Integration in Apache Airflow: Anleitung zur Einbindung der Software in die bestehende Airflow-Instanz.
    \item Funktionsweise der Software: Detaillierte Beschreibung der Anwendung und ihrer Hauptfunktionen.
    \item Fehlermeldungen und Maßnahmen: Übersicht der möglichen Fehlermeldungen und empfohlene Vorgehensweisen zur Fehlerbehebung.
    \item Erweiterungsmöglichkeiten: Anleitung zur Anpassung der Software bei Änderungen der API des Versanddienstleisters, 
    wie z.B. neue Status-Codes oder geänderte CSV-Schemata.
\end{itemize}

\subsection{Code Dokumentation}
\label{sec:CodeDokumentation}

Der Programmcode enthält umfassende Python-Docstrings, die die Funktionen und Klassen erläutern. 
Diese richten sich an Personen, die den erstellten Programmcode anpassen oder besser verstehen möchten. 
Ergänzend dazu finden sich im Anhang:
\begin{itemize}
    \item A9: Entwicklerdokumentation (Auszug Wiki)
    \item A10: Entwicklerdokumentation (Auszug Git)
    \item A6: Programmcode (Auszug \textit{bc\_db.py})
    \item A7: Programmcode (Auszug \textit{trackings.py})
\end{itemize}

% !TEX root = ../Projektdokumentation.tex
\section{Fazit}
\label{sec:Fazit}

\subsection{Soll-/Ist-Vergleich}
\label{sec:Soll-/Ist-Vergleich}
Das Projektziel, ein automatisiertes System zur Überwachung und Benachrichtigung von Warensendungen zu entwickeln, wurde erfolgreich erreicht.
Während des Projekts kam es jedoch zu zeitlichen Abweichungen von der ursprünglichen Planung. Besonders die Analyse- und Implementierungsphase nahmen mehr Zeit in Anspruch als zunächst kalkuliert. 
Die genauen Gründe für diese Verzögerungen sind der detaillierten Zeitplanung \Anhang{app:Zeitplanung} zu entnehmen. Solche Abweichungen sind im Rahmen komplexer Projekte nicht ungewöhnlich und wurden durch eine flexible Anpassung 
des Projektplans erfolgreich aufgefangen. Insgesamt konnte das Projekt trotz dieser Anpassungen im Rahmen der geplanten Ressourcen abgeschlossen werden. 
Die gewonnenen Erkenntnisse im Umgang mit unvorhergesehenem Mehraufwand und die Fähigkeit zur Anpassung der Planung sind wertvolle Erfahrungen für zukünftige Projekte.



% Literatur ------------------------------------------------------------------
\clearpage
\renewcommand{\refname}{Literaturverzeichnis}
\bibliography{Bibliographie}
\bibliographystyle{Allgemein/natdin} % DIN-Stil des Literaturverzeichnisses


% Anhang ---------------------------------------------------------------------
\clearpage
\appendix
\pagenumbering{roman}
% !TEX root = Projektdokumentation.tex
\section{Anhang}

\subsection{Lastenheft}
\label{app:Lastenheft}
\begin{figure}[htb]
\centering
\includegraphicsKeepAspectRatio{Lastenheft_snipped.PNG}{.8}
\caption{Auszug Lastenheft}
\end{figure}
\clearpage
\clearpage

\subsection{Komplette Zeitplanung}
\label{app:ZeitplanungKomplett}

\begin{tabularx}{\textwidth}{Xrrr}
\label{app:Zeitplanung}\\
\rowcolor{heading}\textbf{Analysephase} & \textbf{} & \textbf{} & \textbf{10 h} \\
1. Erhebung der Anforderungen für die Speicherung und Verarbeitung der Trackingdaten &       &    & 3 h  \\
\rowcolor{odd}2. Analyse der bestehenden Systeme und Identifikation von Schwachstellen &       &    & 4 h  \\
3. Ermittlung der Projektkosten und Berechnung der Einsparungen &       &    & 2 h \\
\rowcolor{odd}4. Erstellung des Wirtschaftlichkeitsberichts &       &    & 1 h  \\
\rowcolor{heading}\textbf{Entwurfsphase} & \textbf{} & \textbf{} & \textbf{13 h} \\
1. Entwurf der Architektur für die Datenbank und Integration mit dem ERP-System (Architekturdiagramm) &       &    & 5 h \\
\rowcolor{odd}2. Design der Datenbankstruktur und der Schnittstellen für die Datenübertragung &       &    & 5 h \\
3. Anbindung an das Monitoring- und Benachrichtigungssystem &       & 3 h   &  \\
\rowcolor{heading}\textbf{Implementierungsphase} & \textbf{} & \textbf{} & \textbf{38 h} \\
1. Implementierung der Schnittstellen zur Abholung und Speicherung der Trackingdaten in der Datenbank &       &    & 10 h  \\
\rowcolor{odd}2. Entwicklung der Integration des ERP-Systems zur Synchronisation und Datenabgleich &       &    & 10 h  \\
3. Programmierung der Prüfalgorithmen für die Trackingdaten und Einrichtung des Monitoring-Systems mit Python &       &    & 8 h  \\
\rowcolor{odd}4. Implementierung des Benachrichtigungssystems für die Fachabteilung &       &    & 4 h \\
5. Integration aller Komponenten und Durchführung erster Funktionstests &       &    & 3 h \\
\rowcolor{odd}6. Debugging und Behebung von initialen Fehlern &       &    & 3 h \\
\rowcolor{heading}\textbf{Testphase} & \textbf{} & \textbf{} & \textbf{8 h} \\
1. Durchführung von Funktionstests zur Sicherstellung der korrekten Datenverarbeitung und -integration &       &    & 4 h \\
\rowcolor{odd}2. Validierung der Datenintegrität und Konsistenzprüfungen &       &    & 2 h \\
3. Testen des Monitoring-Systems und der Benachrichtigungsfunktionalität &       &    & 2 h \\
\rowcolor{heading}\textbf{Einführungsphase} & \textbf{} & \textbf{} & \textbf{3 h} \\
1. Deployment des Systems auf die Produktionsumgebung und Durchführung der finalen Systemkonfiguration &       &    & 2 h \\
\rowcolor{odd}2. Schulung der Fachabteilung zur Nutzung des neuen Systems &       &    & 1 h \\
\rowcolor{heading}\textbf{Dokumentation} & \textbf{} & \textbf{} & \textbf{8 h} \\
1. Erstellung der technischen Dokumentation, einschließlich Systemarchitektur und Datenflussdiagramme &       &    & 8 h \\
\hline
\hline
\rowcolor{heading}\textbf{Gesamt} & \textbf{} & \textbf{} & \textbf{80 h} \\
\end{tabularx}


\subsection{Zeitplanung nach Anpassungen}
\begin{tabularx}{\textwidth}{Xrrr}
\label{app:Zeitnachher} \\
\rowcolor{heading}\textbf{Phase} & \textbf{Geplant} & \textbf{Tatsächlich} & \textbf{Differenz} \\
\textbf{Analysephase} & 10 h   & 12 h  & 2 h \\
\rowcolor{odd}\textbf{Entwurfsphase} & 13 h  & 12 h  & -1 h \\
\textbf{Implementierungsphase} & 38 h  & 38 h  &   \\
\rowcolor{odd}\textbf{Testphase} & 8 h   & 6 h   & -2 h  \\
\textbf{Einführungsphase} & 3 h   & 3 h   &  \\
\rowcolor{odd}\textbf{Dokumentation Erstellen} & 8 h   & 9 h  & 1 h \\
\hline
\hline
\rowcolor{heading}\textbf{Gesamt} & 80 h  & 80 h  &  \\
\end{tabularx}


\subsection{PostgreSQL}
\label{app:PostgreSQL}
\begin{figure}[htb]
\centering
\includegraphicsKeepAspectRatio{sql.png}{1}
\caption{DBeaver - SQL Editor}
\end{figure}


\subsection{Airflow}
\label{app:Airflow}
\begin{figure}[htb]
\centering
\includegraphicsKeepAspectRatio{Airflow.png}{1}
\caption{Airflow: Übersicht der Benutzeroberfläche}
\end{figure}


\subsection{Abnahmephase}
\label{app:Abnahmephase}
\begin{figure}[htb]
\centering
\includegraphicsKeepAspectRatio{dummy_daten_sftp.png}{.8}
\caption{Filezilla: Dummy-Daten auf dem SFTP-Server}
\end{figure}
\clearpage









\end{document}


